%======================================================================
% Section 1: Review Scope and Purpose
%======================================================================
\section{Review Scope and Purpose}
\label{sec:scope}

\subsection{Motivation and scientific context}

Wave propagation in generalized continua sits at the intersection of classical continuum mechanics, thermodynamics, and coupled-field physics, and the field exists because of a single contradiction that the classical theory cannot resolve. Fourier's law, together with the balance of energy, implies that a temperature disturbance is felt everywhere in the body at the same instant. Biot's synthesis of that law with linear elasticity \cite{Biot1956} was, and remains, the most successful continuum model of thermal-mechanical interaction ever constructed; it is accurate, economical, and it carries with it an infinite propagation speed that is simply wrong. Every subsequent attempt to repair it has been a repair of this one defect, and the resulting sixty-year history is best understood not as a sequence of competing theories but as a systematic exploration of the ways in which a constitutive law can remember.

Those ways turn out to be few, and they recur. A constitutive law can be given an internal variable that relaxes towards its equilibrium value, which produces the relaxation-time family --- a single time scale attached to the heat flux, as in the Lord--Shulman construction \cite{Lord1967}, or two, as in the two-temperature Green--Lindsay theory \cite{Green1972}, or a thermal displacement that removes the dissipation altogether in the conservative Green--Naghdi construction \cite{Green1991,Green1993}. It can be given separate lags for the flux and for the temperature gradient, which is the phase-lag construction of Tzou \cite{Tzou1995,Tzou1997} and the sharpest statement that the thermal and mechanical laws may be generalised at different rates. It can be given an entire memory kernel, which is the hereditary and memory-dependent construction of Wang and Li \cite{Wang2011}. It can be given a fractional order, which replaces the discrete choice between diffusion and conduction with a continuum of intermediate behaviour \cite{Caputo1967,Povstenko2004,Povstenko2016}. Or the generalising parameter can be put into space rather than into time, as in Eringen's nonlocal elasticity \cite{Eringen1972,Eringen1974,Eringen1983}, in which the stress at a point is an average over a neighbourhood and the material therefore acquires a length scale of its own.

Each of these repairs removes the infinite propagation speed and imports a new cost. The second temperature must be supplied with initial data. The thermal displacement changes the character of the energy, and with it the admissibility conditions. The phase lag introduces a parameter that must satisfy an ordering condition for the second law to hold. The memory kernel makes the constitutive relation nonlocal in time and the discretisation history-dependent. The fractional order makes the initial-value problem well posed only in a suitable function space. The nonlocal length makes the order of the governing equation higher than the classical one, and the local well-posedness of the initial-value problem is no longer automatic. These costs are the substance of the field: a review that reported only the benefits would misrepresent it, and a review that reported only the costs would miss that the benefits are real and, in their regimes, decisive.

Alongside this thermal axis runs a second one. A material that deforms under load generally does other things as well: it develops an electrical polarization, it changes its void content, it rotates its internal microstructure, it generates and moves carriers. Adding any one of these to the thermoelastic framework produces a coupled multiphysics continuum --- piezothermoelasticity, thermoporous and double-porous media, micropolar elasticity, semiconductor photo-thermo-elasticity, magneto-thermo-elasticity --- and the addition is not merely an extra term. It adds a field, and therefore adds a mode, and therefore enlarges the modal structure that the wave analysis must resolve, and adds a matching condition, and therefore enlarges the interface system that the scattering analysis must solve. The simultaneous presence of both axes --- memory in time and coupling in space --- is where the difficulty of the field is concentrated, and where the research questions of Section~\ref{sec:future} are located.

These coupled formulations underpin smart structures, energy harvesters, nondestructive evaluation, geophysical prospecting, biomedical ultrasound and nanodevice design, and they have produced a substantial literature in their own right. Yet the field lacks the kind of synthesis that a practitioner needs: a single account that states the constitutive structure of each family, derives the wave problem it produces, establishes the conditions under which each is thermodynamically admissible, shows how the resulting quantities are computed, and identifies what remains open. Existing reviews tend to cover one family or one application \cite{Shakeriaski2021,Chandrasekharaiah1998,Joseph1989,Askes2006}; this review is organised so that the same question can be followed across all of them.

\subsection{Objectives and contributions}

This review has the following objectives:

\begin{enumerate}
\item To present the first-principles foundations of continuum wave mechanics: kinematics, balance laws, thermodynamic restrictions, and constitutive principles (Section 3).
\item To trace the historical development of generalized thermoelastic theories from classical coupled thermoelasticity to modern memory-dependent and nonlocal formulations, explaining the physical motivation behind each theory (Section 4).
\item To present coupled multiphysics continua—piezothermoelasticity, porous media, micropolar continua, semiconductor interactions, and magneto-thermo-elasticity—with their governing equations and coupling mechanisms (Section 5).
\item To survey generalized transport and memory effects: relaxation, fractional operators, hereditary integrals, memory-dependent derivatives, and their mathematical structure (Section 6).
\item To present nonlocal and scale-dependent continuum models, including Eringen's nonlocal elasticity and strain-gradient theory, with their wave dispersion characteristics (Section 7).
\item To review mathematical methods for wave propagation analysis: plane wave ansatz, Christoffel equations, secular equations, transform methods, and numerical approaches (Section 8).
\item To develop the theory of wave reflection, transmission, and energy partitioning at interfaces between different continuum media (Section 9).
\item To discuss energy transport and dissipation across theory families, including the verified piezothermoelastic energy flux and open questions for pure generalized theories (Section 10).
\item To survey advanced theoretical developments and emerging research directions (Sections 11--12).
\item To present applications with evidence levels, an original five-axis research taxonomy, major scientific challenges, and future research opportunities (Sections 13--16).
\end{enumerate}

\subsection{Scope and organization}

The review covers the following theory families:

\begin{itemize}
\item \textbf{Classical coupled thermoelasticity} (Biot, 1956) — Fourier heat conduction coupled with elastic wave propagation.
\item \textbf{Lord--Shulman theory} (1967) — generalized thermoelasticity with one thermal relaxation time.
\item \textbf{Green--Lindsay theory} (1972) — two-temperature thermoelasticity with two relaxation times.
\item \textbf{Green--Naghdi types I--III} (1991, 1993) — thermal displacement formulation, including dissipationless type II.
\item \textbf{Dual-phase-lag and triple-phase-lag theories} (Tzou, 1995/1997) — phase-lag heat conduction models.
\item \textbf{Moore--Gibson--Thompson theory} — third-order heat conduction with thermal inertia.
\item \textbf{Memory-dependent derivative models} (Wang \& Li, 2011) — integral-form memory operators with fading memory.
\item \textbf{Fractional thermoelasticity} — time-fractional heat equations for anomalous diffusion.
\item \textbf{Nonlocal elasticity} (Eringen, 1972--1983) — scale-dependent constitutive relations with spatial memory.
\item \textbf{Piezothermoelasticity and coupled multiphysics continua} — piezoelectric, porous, micropolar, semiconductor, and magneto coupling.
\item \textbf{Interface wave phenomena} — reflection, transmission, scattering, and energy partitioning at material interfaces.
\end{itemize}

The manuscript is organized into seventeen sections:

\begin{itemize}
\item \textbf{Section 1} (this section): scope, purpose, and organization.
\item \textbf{Section 2}: literature landscape and scholarly context, including comparison with existing reviews.
\item \textbf{Sections 3--5}: Layer I--III of the mathematical exposition—foundations, thermoelastic theories, and coupled multiphysics continua.
\item \textbf{Sections 6--8}: generalized transport, memory effects, nonlocal models, and mathematical methods.
\item \textbf{Sections 9--10}: interface phenomena and energy transport.
\item \textbf{Sections 11--12}: advanced developments and comparative synthesis.
\item \textbf{Sections 13--17}: applications, taxonomy, challenges, future directions, and conclusions.
\end{itemize}

Appendices provide an equation verification register, balance law derivations, numerical benchmarks, figure sources, table sources, and reproducibility notes.

\subsection{Intended readership}

This review is written for multiple audiences:

\begin{itemize}
\item \textbf{Beginning researchers and MSc students:} Sections 1--3 provide the foundational balance laws, constitutive relations, and classical wave mechanics needed to enter the field. The notation tables and Section 17 provide orientation.
\item \textbf{Doctoral researchers:} Sections 3--10 provide mathematical depth with representative equations, dispersion relations, interface formulations, and solution methods. Section 16 identifies dissertation topics.
\item \textbf{Experienced scientists:} Sections 12--16 provide comparative synthesis, taxonomy, challenge analysis, and future research roadmap.
\item \textbf{Journal reviewers and editors:} Sections 1--2, 12, and 16--17 provide context for evaluating new submissions in generalized continuum mechanics.
\end{itemize}

\subsection{Scholarly principles}

This review adheres to the following principles:

\begin{enumerate}
\item \textbf{Constructive synthesis:} Model-specific differences are traced to documented assumptions (constitutive, material symmetry, geometry, boundary conditions, thermal theory, frequency convention, normalization, approximation level), not presented as contradictions or errors.
\item \textbf{Mathematical accuracy:} Equations are drawn from verified derivations or established primary sources. Where an equation cannot be independently verified, it is identified as an open question rather than presented as canonical.
\item \textbf{Comparative rather than adjudicative framing:} Competing formulations are compared on the basis of their modelling choices and domains of validity; the review does not attempt to adjudicate between individual publications.
\item \textbf{Minimal equation set:} Fewer, precise equations are preferred over hundreds of redundant ones. Representative equations are selected to illustrate each theory family.
\item \textbf{Original figures and tables:} All figures are original vector graphics created for this review. All tables synthesize knowledge across theory families rather than listing individual papers.
\end{enumerate}


\begin{figure}[htbp]
\centering
\includegraphics[width=0.9\textwidth]{fig01-theory-evolution.pdf}
\caption{Historical evolution of thermoelasticity theories from Fourier's heat equation (1822) to modern phase-lag and memory-dependent formulations (1995+). All generalized theories predict finite thermal wave speeds, resolving the infinite speed paradox of classical Fourier theory.}
\label{fig:theory_evolution}
\end{figure}

\begin{figure}[htbp]
\centering
\includegraphics[width=0.9\textwidth]{fig02-theory-relationships.pdf}
\caption{Relationships and limiting cases between thermoelastic theory families. All generalized theories converge to classical coupled thermoelasticity as relaxation parameters approach zero. These are parameter substitutions, not rigorous limit theorems.}
\label{fig:theory_relationships}
\end{figure}

% LocalWords: 
